%% AMS-LaTeX Created with the Wolfram Language : www.wolfram.com

\documentclass[12pt]{article}

\usepackage[letterpaper,margin=1in]{geometry}
\usepackage{amsmath, amssymb, graphics, setspace}
\usepackage{graphicx} 
\graphicspath{ {./figs/} }
\usepackage{tcolorbox}
\usepackage{bm}
\usepackage[framed,numbered,autolinebreaks,useliterate]{mcode}
\newcommand{\mathsym}[1]{{}}
\newcommand{\unicode}[1]{{}}

\usepackage{hyperref}
\hypersetup{
    colorlinks=true,
    linkcolor=blue,
    filecolor=magenta,      
    urlcolor=cyan,
    pdftitle={Overleaf Example},
    pdfpagemode=FullScreen,
}

\def\mathbi#1{\textbf{\em #1}}

\begin{document}
\doublespacing	

\title{CE 401/5501 Homework 11\\
\vspace{2in}
\textbf{Name: \rule{2in}{0.15mm}}}
\author{}
\maketitle

\newcommand*{\I}{\imath} % imaginary unit i

\newpage
\doublespacing

\section*{Overview}
In this assignment, we will extend our exploration of the MNIST dataset by:
\begin{enumerate}
  \item Evaluating the impact of dimensionality reduction (PCA and t-SNE) on classification performance using a Perceptron.
  \item Investigating how the regularization parameter \(C\) affects multi‐class SVM (with RBF kernel) performance, and selecting the best \(C\) via cross‐validation.
  \item Put it into action: experiment and test your own capture!
\end{enumerate}

In the meantime, please review 
\begin{enumerate}
    \item  The \href{https://umsystem.instructure.com/courses/280496/files/33966689?module_item_id=9952773}{class presentation} about the NISt/MNIST dataset. 
    \item The \href{https://umsystem.instructure.com/courses/280496/files/33966689?module_item_id=9952773}{MNIST notebook} is here and you should copy it to your google's colab notebook folder to run it. 
    \item Dimensionality demonstration as in this \href{https://umsystem.instructure.com/courses/280496/files/33859915?module_item_id=9938963}{notebook}.
\end{enumerate}

\newpage

\section*{Task 1: Dimensionality Reduction \& Perceptron}

Please following the following steps to complete:
\subsection*{A. Data Preparation}
\begin{enumerate}
  \item Load MNIST from \texttt{sklearn.datasets}.
  \item Standardize the pixel values to zero mean and unit variance.
  \begin{itemize}
      \item We didn't do this in class as we used the `mnist\_784' dataset directly. 
      \item % Given a pixel vector \bm{x} = (x_1, \dots, x_d)^\top, we compute its sample mean and standard deviation, then standardize each component.

\[
\text{Mean:}\quad
\mu = \frac{1}{d} \sum_{i=1}^d x_i,
\qquad
\text{(the average pixel value)}
\]

\[
\text{Standard deviation:}\quad
\sigma = \sqrt{\frac{1}{d} \sum_{i=1}^d \bigl(x_i - \mu\bigr)^2},
\qquad
\text{(the spread of the pixel values)}
\]

\[
\text{Standardization:}\quad
\tilde{x}_i = \frac{x_i - \mu}{\sigma}, 
\quad i = 1,2,\dots,d,
\]
\[
\text{so that the transformed vector }
\tilde{\bm{x}} = (\tilde{x}_1,\dots,\tilde{x}_d)^\top
\text{ has zero mean and unit variance.}
\]

  \end{itemize}
\end{enumerate}

If programmed properly, the entire dataset of $\bm{X}$ can be processed and standardized in a few lines. Please explore this. We will use the symbol $\bm{Y}$ to express it.

\subsection*{B. PCA (Principal Component Analysis)}
\begin{enumerate}
  \item Use \texttt{sklearn.decomposition.PCA} to reduce the 784‐dimensional standardized data $\bm{Y}$ to 3 principal components.
  \item Create a 3D scatter plot or three 2D scatter plots of the first 3 components, coloring points by their true digit labels. 
  \item Report the scatter plots; and remark that if you see 10 clusters well separated corresponding to the underlying 10 digits. 
\end{enumerate}

\subsection*{C. t‑SNE}
\begin{enumerate}
  \item Use \texttt{sklearn.manifold.TSNE} with \texttt{n\_components=3}, \texttt{perplexity=30}, and a select random state.
  \item Create a 3D scatter plot or three 2D scatter plots of the first 3 components, coloring points by their true digit labels. 
  \item Report the scatter plots; and remark that if you see 10 clusters well separated corresponding to the underlying 10 digits. 
  \item Comment on visual separability compared to PCA.
\end{enumerate}

\subsection*{D. Classification with Perceptron}
\begin{enumerate}
  \item Train a \texttt{sklearn.linear\_model.Perceptron} on:
    \begin{itemize}
      \item The original 784‐dimensional data.
      \item The 3‐dimensional PCA‐reduced data.
      \item The 3‐dimensional t‑SNE‐reduced data.
    \end{itemize}
  \item Evaluate each model on a held‐out test set (i.e., the same testing data); report accuracy, precision, recall, and F1‐score for each digit. 
  \item Summarize how dimensionality reduction affected classification performance.
\end{enumerate}

\newpage

\section*{Task 2: SVM Hyperparameter Tuning}
\subsection*{Baseline SVM}
\begin{enumerate}
  \item Using the full (no feature‐extraction) MNIST dataset, train a multi‐class SVM with RBF kernel (\texttt{sklearn.svm.SVC(decision\_function\_shape='ovr')}). \textit{(note that here we are using `one-vs-rest' multi-class scheme).}
  \item Use default \(C=1.0\) as a baseline; report the test set accuracy.
\end{enumerate}

\subsection*{Effect of the \(C\) Parameter}
\begin{enumerate}
  \item Define a grid of \(C\) values, e.g.\ \(\{0.01, 0.1, 1, 10, 100\}\).
  \item For each \(C\), perform 5‐fold cross‐validation on a fraction of the training set (you can choose $20\%$ of the total training dataset), and record the mean CV accuracy. 
  \begin{itemize}
      \item K-fold cross-validation is an essential technique for hyperparameter selection or more generally model selection toward generalization. 
      \item Learn the usage of `cross-validation' from \href{https://scikit-learn.org/stable/modules/cross_validation.html#multimetric-cross-validation}{this link}.
  \end{itemize}
  \item Plot \(\mathrm{CV\ Accuracy}\) vs.\ \(\log_{10}(C)\), and report the resulting plot. 
\end{enumerate}

\subsection*{Selecting the Best \(C\)}
\begin{enumerate}
  \item Identify the \(C\) value with the highest mean cross‐validation accuracy.
  \item Retrain the SVM using this \(C\) on the full training set.
  \item Evaluate on the test set; report the final test accuracy and a confusion matrix (\texttt{sklearn.metrics.ConfusionMatrixDisplay}).
\end{enumerate}

\section*{Task 3: Test your own image}
Please review the code added at the last of the class demo \href{https://umsystem.instructure.com/courses/280496/files/33967083?module_item_id=9952773}{notebook}

\begin{itemize}
    \item Try the best calligraphic skills and write down some numbers on paper (each number on one paper), and take pictures using phone. Then transfer it to your google drive. 
    \item Learn the code there and convert it to MNIST ready feature vectors and call your best hand-written digit classification model so far and test the accuracy.
    \item You may check out \href{https://umsystem.instructure.com/courses/280496/files/33966714?module_item_id=9952778}{this image} - the logistic regression classifier used in class actually failed!
    \item Report your results and remark. 
\end{itemize}

\end{document}
